\scanpage{63}
\margin{\rotatebox{180}{62}}
\mspage{- 7 -}

$(B_q)$. \uline{Il existe un sous-groupe invariant} $C$ \uline{de} $G$
\uline{qui est cyclique d'ordre premier à} $q$, \uline{et qui est tel que}
$G/C$ \uline{soit un} $q$-\uline{groupe}.

Le résultat suivant est dû à Brauer [3] et Witt [18]:

PROPOSITION 2. \uline{Soit} $G$ \uline{un groupe fini et soit} $E$
\uline{un corps de caractéristique zéro}. \uline{Supposons que, pour tout nombre premier}
$q$, \uline{et pour tout sous-groupe} $H$ \uline{de} $G$ \uline{vérifiant}
$(B_q)$, \uline{l'algèbre} $E[H]$ \uline{soit décomposée}. \uline{Alors}
$E[G]$ \uline{est décomposée}.

Rappelons brièvement la démonstration:

On commence par établir une formule:
\[
(12)\qquad 1=\underset{\lambda}{\quad}n_\lambda\varphi^*_\lambda,
\quad n_\lambda\in\underline{\underline Z},
\]
où les $\varphi_\lambda$ sont des caractères de représentations rationnelles sur
$E$ de sous-groupes $H_\lambda$ vérifiant $(B_{q_\lambda})$, et où
$\varphi^*_\lambda$ désigne le caractère de $G$ \uline{induit} par
$\varphi_\lambda$; voir Witt [18], Satz 6, ou bien Swan [14], th.4.1.

Si maintenant $\chi$ est un caractère de $G$ à valeurs dans $E$, on multiplie
(12) par $\chi$, et l'on obtient:
\[
(13)\qquad \chi=\sum n_\lambda\varphi^*_\lambda\chi
=\sum n_\lambda(\varphi_\lambda\chi)^*.
\]
Comme $\varphi_\lambda\chi$ est un caractère de $H$ à valeurs dans $E$,
l'hypothèse faite sur les $H_\lambda$ montre que c'est le caractère d'une
représentation de $H_\lambda$ rationnelle sur $E$; donc
$(\varphi_\lambda\chi)^*$ est le caractère d'une représentation de $G$
rationnelle sur $E$. La formule (13), jointe au critère (b) ci-dessus, montre
que $\chi$ est le caractère d'une représentation de $G$ rationnelle sur $E$.
Ceci ayant lieu pour tout $\chi$, la prop.1 montre que $E[G]$ est décomposée, cqfd.

3. \uline{Démonstration du théorème 3}.

Compte tenu de la prop.1, il s'agit de prouver que, \uline{si} $G$ \uline{vérifie}
$(R_p)$ \uline{et si} $\ell\ne p$, \uline{l'algèbre} $Q_\ell[G]$
\uline{est décomposée}. De plus, la prop.2 permet de supposer qu'il existe un
nombre premier $q$ tel que $G$ vérifie $(B_q)$. Nous
